\section{Resumen del capítulo}

\begin{tcolorbox}[resumen,title=Producto]
El producto es una operación binaria. En los números naturales puede interpretarse como una suma repetida:
\[
a\cdot b=\underbrace{b+b+\cdots+b}_{a\text{ veces}}.
\]
Los números que participan reciben el nombre de \emph{factores}. Cuando una expresión combina sumas y productos, el producto se agrupa antes que la suma.
\end{tcolorbox}

\begin{tcolorbox}[resumen,title=Propiedades del producto]
Para los números naturales, enteros y racionales, el producto cumple:
\begin{gather*}
a\cdot b=b\cdot a\\
(a\cdot b)\cdot c=a\cdot(b\cdot c)\\
a\cdot1=a\\
a\cdot0=0\\
a\cdot(b+c)=a\cdot b+a\cdot c
\end{gather*}
Estas propiedades permiten reorganizar y transformar productos sin modificar el número que representan.
\end{tcolorbox}

\begin{tcolorbox}[resumen,title=Producto de enteros]
El producto se extiende de los naturales a los enteros conservando sus propiedades. El producto por \(-1\) produce el opuesto:
\[
a\cdot(-1)=-a.
\]
Por ello, el producto de factores con el mismo signo es positivo y el de factores con distinto signo es negativo; todo producto por cero es cero.
\end{tcolorbox}

\begin{tcolorbox}[resumen,title=Números racionales]
Los números racionales son los que pueden escribirse como
\[
  \frac{a}{b}
\]
con \(a,b\in\mathbb Z\) y \(b\ne0\).

Una misma cantidad puede tener fracciones equivalentes:
\[
\frac{a}{b}=\frac{ac}{bc},\qquad c\ne0.
\]
Además,
\[
\frac{a}{b}\cdot\frac{c}{d}=\frac{a\cdot c}{b\cdot d}
\]
\end{tcolorbox}

\begin{tcolorbox}[resumen,title=Inverso y división]
Todo racional no nulo tiene un inverso multiplicativo:
\[
  \left(\frac{a}{b}\right)^{-1}=\frac{b}{a}\quad \text{con } a\ne0.
\]
La división por un número no nulo se define mediante ese inverso:
\[
x\div y=x\cdot y^{-1}.
\]
El cero no tiene inverso multiplicativo.
\end{tcolorbox}

\begin{tcolorbox}[resumen,title=Representaciones y conjuntos numéricos]
Una fracción, una división y una escritura decimal pueden representar el mismo número; por ejemplo,
\[
1\div2=\frac12=0{,}5.
\]
La escritura decimal no constituye un conjunto numérico nuevo. Los conjuntos estudiados se incluyen según
\[
\mathbb N\subseteq\mathbb Z\subseteq\mathbb Q.
\]
\end{tcolorbox}
